% 奥利弗·亥维赛（综述）
% license CCBYSA3
% type Wiki

本文根据 CC-BY-SA 协议转载翻译自维基百科\href{https://en.wikipedia.org/wiki/Oliver_Heaviside}{相关文章}

奥利弗·赫维赛德（Oliver Heaviside，/ˈhɛvisaɪd/，读作“赫维赛德”\(^\text{[2]}\)，1850年5月18日－1925年2月3日）是一位英国数学家和电气工程师。他发明了一种求解微分方程的新方法（与拉普拉斯变换等价），独立发展了向量分析，并将麦克斯韦方程组改写成如今通用的形式。他对麦克斯韦方程组在麦克斯韦逝世后的理解与应用产生了深远影响。此外，1893年，赫维赛德将麦克斯韦方程扩展到了“引力电磁学”领域——这一理论在2005年的“重力探测器B”实验中得到了验证。他提出的“电报方程”在他生前的晚期获得了商业上的重要意义，尽管这一成果曾长期未被重视，因为当时很少有人理解他所采用的创新数学方法。\(^\text{[3]}\)尽管赫维赛德一生大部分时间都与主流科学界格格不入，但他彻底改变了电信、数学和科学的面貌。\(^\text{[3]}\)
